\chapter{ROS Workspace Architecture}

\section{Repository layout}

The robot repository is the workspace at:

\begin{lstlisting}[language=bash]
/home/rudra/Projects/RUDRAv4
\end{lstlisting}

The primary package is:

\begin{lstlisting}[language=bash]
src/rudra_base_bridge
\end{lstlisting}

The package is Python-based, but it carries firmware sketches, YAML configuration, RViz files, and launch files. This is acceptable for a robot-development workspace because the package is the unit of bringup, not merely the unit of Python code. The normal operator launch is \filep{rudra_full.launch.py}; lower-level launch files and shell wrappers remain available for commissioning individual subsystems.

\section{Node inventory}

\begin{longtable}{p{0.25\linewidth}p{0.25\linewidth}p{0.40\linewidth}}
\toprule
Executable & Source file & Purpose\\
\midrule
\pkg{ps2_uno_to_teensy} & \filep{ps2_uno_to_teensy.py} & Reads \cmd{J,...} from Uno, applies joystick mixing and obstacle guard, writes \cmd{D,...} to Teensy, publishes observability topics.\\
\pkg{cmd_vel_to_teensy} & \filep{cmd_vel_to_teensy.py} & Converts ROS \topic{/cmd_vel} into left/right Sabertooth commands and writes them to Teensy.\\
\pkg{dcdc_usb_monitor} & \filep{dcdc_usb_monitor.py} & Parses \cmd{dcdc-usb -a} output and publishes NUC battery state and diagnostics.\\
\pkg{odom_tf_broadcaster} & \filep{odom_tf_broadcaster.py} & Broadcasts \cmd{odom -> base_link} from odometry messages when localization is active.\\
\pkg{serial_port_list} & \filep{serial_port_list.py} & Helper to list serial ports.\\
\bottomrule
\end{longtable}

\section{Helper modules}

Three helper modules are central to maintainability:

\begin{description}
\item[\pkg{serial_utils.py}] Provides \pkg{LineSerial}, a reconnecting line-oriented wrapper around pyserial, plus clamp, map, and deadband helpers.
\item[\pkg{obstacle_guard.py}] Encapsulates LiDAR scan filtering and command scaling. Both bridge nodes use the same guard behavior.
\item[\pkg{teensy_telemetry.py}] Parses optional \cmd{IMU,...} and \cmd{ODOM,...} lines and publishes ROS messages.
\end{description}

This modularity matters because the two drive bridges should not maintain two independent obstacle-guard implementations or two independent telemetry parsers.

\section{Launch file inventory}

\begin{longtable}{p{0.31\linewidth}p{0.59\linewidth}}
\toprule
Launch file & Role\\
\midrule
\filep{rudra_full.launch.py} & Main operator launch. Starts PS2 manual bridge, Teensy serial owner, LiDAR, obstacle guard, localization, voice, and optional Ollama, RViz, and DCDC monitoring.\\
\filep{rudra_ps2.launch.py} & Starts manual PS2 bridge only.\\
\filep{rudra_ps2_lidar.launch.py} & Starts PS2 bridge, YDLIDAR driver, and static \cmd{base_link -> laser} transform.\\
\filep{rudra_ps2_lidar_localization.launch.py} & Includes PS2+LiDAR launch, includes localization launch, and can optionally start DCDC monitor.\\
\filep{cmd_vel_to_teensy.launch.py} & Starts the autonomous or teleop \topic{/cmd_vel} bridge.\\
\filep{lidar.launch.py} & Starts only the YDLIDAR driver and static LiDAR transform.\\
\filep{lidar_view.launch.py} & Starts RViz view for scan and obstacle marker. It does not start a driver.\\
\filep{localization.launch.py} & Starts EKF, odom TF broadcaster, and \cmd{base_link -> imu_link} static transform.\\
\filep{localization_view.launch.py} & Starts RViz view for odometry, fused odometry, scan, and transforms.\\
\filep{dcdc_usb_monitor.launch.py} & Starts NUC power monitor only.\\
\filep{rudra_voice.launch.py} & Starts the voice node and command guard. The full launch includes this when \cmd{enable_voice:=true}.\\
\bottomrule
\end{longtable}

\begin{figure}[htbp]
\centering
\resizebox{0.98\linewidth}{!}{\begin{tikzpicture}[node distance=7mm]
\node[rudraBlock] (full) {\filep{rudra_full.launch.py}};
\node[rudraBlock,below left=of full] (base) {\filep{rudra_ps2_lidar_localization.launch.py}};
\node[rudraBlock,below right=of full] (voice) {\filep{rudra_voice.launch.py}};
\node[rudraBlock,below=of full] (helpers) {optional\\Ollama, RViz, DCDC};
\node[rudraBlock,below left=of base] (ps2lidar) {\filep{rudra_ps2_lidar.launch.py}};
\node[rudraBlock,below right=of base] (loc) {\filep{localization.launch.py}};
\node[rudraTopic,below left=of ps2lidar] (ps2) {\pkg{ps2_uno_to_teensy}};
\node[rudraTopic,below=of ps2lidar] (lidar) {\pkg{ydlidar_ros2_driver_node}};
\node[rudraTopic,below right=of ps2lidar] (tf1) {static TF\\\cmd{base_link -> laser}};
\node[rudraTopic,below left=of loc] (ekf) {\pkg{ekf_filter_node}};
\node[rudraTopic,below=of loc] (tf2) {\pkg{odom_tf_broadcaster}};
\node[rudraTopic,below right=of loc] (imu_tf) {static TF\\\cmd{base_link -> imu_link}};
\draw[rudraArrow] (full) -- (base);
\draw[rudraArrow] (full) -- (voice);
\draw[rudraDashed] (full) -- (helpers);
\draw[rudraArrow] (base) -- (ps2lidar);
\draw[rudraArrow] (base) -- (loc);
\draw[rudraArrow] (ps2lidar) -- (ps2);
\draw[rudraArrow] (ps2lidar) -- (lidar);
\draw[rudraArrow] (ps2lidar) -- (tf1);
\draw[rudraArrow] (loc) -- (ekf);
\draw[rudraArrow] (loc) -- (tf2);
\draw[rudraArrow] (loc) -- (imu_tf);
\end{tikzpicture}}
\caption{Launch dependency tree for the full robot bringup.}
\end{figure}

\section{Configuration files}

The three YAML files form the runtime contract:

\begin{description}
\item[\filep{rudra_v4_hardware.yaml}] Base bridge parameters: ports, baud, timeouts, joystick mapping, maximum speeds, obstacle guard parameters, telemetry topics, and DCDC thresholds.
\item[\filep{ydlidar_g2b.yaml}] YDLIDAR parameters: port, frame, baudrate, sample rate, scan range, frequency, intensity mode, and inversion/reversion flags.
\item[\filep{localization_ekf.yaml}] EKF parameters: two-dimensional mode, frames, wheel odom input, IMU input, selected measurement variables, queue sizes, and timeouts.
\end{description}

\section{Build model}

The repository build uses \cmd{colcon}. The minimal robot package build is:

\begin{lstlisting}[language=bash]
cd /home/rudra/Projects/RUDRAv4
bash scripts/install_deps.sh
source /opt/ros/lyrical/setup.bash
colcon build --packages-select rudra_base_bridge
source install/setup.bash
\end{lstlisting}

The checked-in \filep{scripts/build_ws.sh} wraps this idea. The build output goes into \filep{install/}. The runtime source order must still be respected in every terminal that launches robot nodes.

\section{Why wrapper scripts still exist}

The full launch is the normal operator path. Wrapper scripts are retained for subsystem commissioning because they set domain variables, source overlays, select known-good config files, and reduce typing mistakes while testing one boundary at a time. For example:

\begin{lstlisting}[language=bash]
bash scripts/run_ps2_lidar_localization.sh
\end{lstlisting}

is useful when proving the PS2, LiDAR, and EKF stack before adding voice or optional helpers. Once commissioning passes, return to \filep{rudra_full.launch.py} for normal operation.
